\subsection{Review of Long Division}
We divide polynomials using a process very similar to the one we use to divide numbers. How do we divide $485 \div 21$?

\begin{lpic}{4}{8}
%	\draw[dashed,thick] (\value{cols},0) -- (\value{cols},-8);
	\writethink{-1}{8}
%	\regtext{-.25\textwidth}{-1}{Write:};
%	\regtext{.25\textwidth}{-1}{Think:};
\end{lpic}
The process of divison long goes:
\begin{lpic}{14}{5}
\foreach \y/\l in {0/D,1/M,2/S,3/B} {
  \bigtext{1.5}{-\y-1}{\l};
}
\end{lpic}
Each of the numbers in a division problem has a name. In the example above:
\begin{itemize}
\item The number $21$ (that we're dividing by) is called the \makeblank[1.5in]
\item The number $485$ (that we're dividing into) is called the \makeblank[1.5in]
\item The number $23$ (the result of the division) is called the \makeblank[1.5in]
\item The number $2$ (what was left over) is called the \makeblank[1.5in]
\end{itemize}
We can check the results of our division by multiplying and then adding:
\[
\makeblank[.5in]\times\makeblank[.5in]+\makeblank[.5in]=\makeblank[.5in]
\]